\scanpage{16}
$\pi_0G'$ de $\pi_0H'$ par $R'=\mathrm{Coker}\,(\pi_1 H'\to N)$. On fera \unsure{$\cdot$} attention que $\pi_0G$ est lui-$=$ une extension de \deleted{Coker} $\pi_0H$ par $\mathrm{Coker}\,(\pi_1H\to N)$, son image inverse \inserted{$\pi_0 G\times_{\pi_0 H}\pi_0 H'$} par $\pi_0 H'\to\pi_0H$ est une extension de $\pi_0H'$ par $R=\mathrm{Coker}\,(\pi_1 H\to N)$ \deleted{\unsure{$\simeq$ qui}}, non par $R'$. $R$ est un quotient de $R'$ par
\[
\frac{\mathrm{Im}\,(\pi_1 H\xrightarrow{\varphi}N)}
     {\mathrm{Im}\,(\pi_1 H'\xrightarrow{\varphi'}N)}
\simeq \pi_1H/(\mathrm{Ker}\,\varphi+\mathrm{Im}\,\pi_1(u))
\]
donc on a
\[
1\to\pi_1H/(\mathrm{Ker}\,\varphi+\mathrm{Im}\,\pi_1(u))\to\pi_0G'
\xrightarrow{c}\pi_0G\times_{\pi_0H}\pi_0H'\to1
\]
Cel \deleted{\unsure{montre}} \inserted{met} en évidence de façon précise comment le noyau de \unsure{$a$|$c$} est lié au $\pi_1H$ et $\pi_1u$, c'est un quotient de $\mathrm{Coker}\,\pi_1u$.

$\beta$) \uline{Cas $H$ connexe}. Alors l'ext. $G$ est caractérisée, à isom unique près, par un hom
\[
\deleted{\pi\longrightarrow}\qquad
\pi\overset{\mathrm{def}}{=}\pi_1(H)\xrightarrow{\unsure{\rho|\varphi}}\mathfrak z(N)
\]
et peut se décrire à l'aide du groupe \deleted{\unsure{rev}} universel $\widetilde H$ de $H$ (ext. de $H$ par $\pi$) en prenant son ``image directe'' par $\pi\xrightarrow{\varphi}N$

Alors l'extension \deleted{\unsure{de}} $G'$ de $H'$ par $N$ peut s'obtenir en prenant l'image inverse $H'\times_H\widetilde H$ de \deleted{\unsure{l'}}l'ext. $\widetilde H$ de $H$ par $\pi$ via $H'\to H$, et prenant l'image par $\pi\xrightarrow{\varphi}\mathfrak z(N)\hookrightarrow N$. Il n'y a de pb que si $H'$ \deleted{\unsure{est}} n'est pas connexe \inserted{(\unsure{disons|sinon|décrivons} les ext. de $\pi_0H'=1$ sont évidentes !)} -- p. ex $N$ discret. \deleted{discret (plus gén \unsure{si $G$},} \deleted{Il faut que \unsure{de l'on|prendre l'}}
% archive sheet number: 15
